% Generated draft: requires mathematical and source review.
\documentclass[11pt]{article}
\usepackage[margin=25mm]{geometry}
\usepackage{fontspec}
\setmainfont{Latin Modern Roman}
\usepackage{amsmath,amssymb,mathtools,unicode-math}
\setmathfont{Latin Modern Math}
\usepackage{enumitem,needspace,xurl,hyperref}
\setlength{\parindent}{0pt}
\setlength{\parskip}{5pt}
\setlength{\emergencystretch}{3em}
\newcommand{\N}{\mathbb{N}}
\newcommand{\Z}{\mathbb{Z}}
\newcommand{\Q}{\mathbb{Q}}
\newcommand{\R}{\mathbb{R}}
\newcommand{\C}{\mathbb{C}}
\newcommand{\F}{\mathbb{F}}
\newcommand{\A}{\mathbb{A}}
\newcommand{\PP}{\mathbb{P}}
\newcommand{\E}{\mathbb{E}}
\newcommand{\eps}{\varepsilon}
\newcommand{\dd}{\,\mathrm{d}}
\begin{document}
\section*{UnsolvedMath ID: 30004818}
\textbf{Source problem code:} catalogue entry
\subsection{Definitions and mathematical statement}
A \emph{local knot} in a three-manifold \(M\) is an embedded circle contained in the interior of an embedded three-ball in \(M\). Identifying such a ball with a standard ball gives a corresponding classical knot in \(S^3\).

The supplied question asks whether the genus-saving cobordism phenomenon reported for local knots in non-orientable three-manifolds can also occur in an orientable three-manifold. The record does not specify enough conventions to give an unconditional, uniquely determined formalization. The following gives an explicit interpretation whose conventions require confirmation.

Fix a category of embeddings, for example smooth embeddings or locally flat topological embeddings. For oriented local knots \(K_0,K_1\subset M\), let \(I=[0,1]\) and define
\[
 g_M(K_0,K_1)=\min\left\{g(\Sigma):\begin{array}{l}
 \Sigma\subset M\times I\text{ is a compact, connected, orientable,}\\
 \text{properly embedded surface in the chosen category,}\\
 \partial\Sigma=(-K_0)\times\{0\}\sqcup K_1\times\{1\}
 \text{ as oriented boundary}
 \end{array}\right\},
\]
with value \(+\infty\) if no such surface exists. Here \(-K_0\) denotes reversal of the knot orientation. Given chosen ball identifications, let \(\widehat K_0,\widehat K_1\subset S^3\) be the corresponding oriented classical knots, and define \(g_{S^3}(\widehat K_0,\widehat K_1)\) using the same conventions in \(S^3\times I\).

Under this orientable embedded-surface interpretation, the proposed existence question is whether there are an orientable three-manifold \(M\) and two local knots \(K_0,K_1\subset M\), with corresponding classical knots, such that
\[
 g_M(K_0,K_1)<g_{S^3}(\widehat K_0,\widehat K_1)<+\infty.
\]
The allowed class of three-manifolds, the embedding category, the orientation conventions, and the admissible cobordisms are unresolved in the supplied record. Thus this inequality is an explicit candidate interpretation, not an independently established equivalent formulation of the source question. In particular, concordance ordinarily means an annular cobordism, whereas positive-genus cobordisms are also suggested by the source wording; the intended restriction is not supplied.
\subsection{Short English statement}
Can an orientable three-manifold provide a cobordism between two knots contained in balls whose genus is smaller than the least possible genus between their corresponding classical knots in \(S^3\)? This expresses a candidate reading of the supplied question. The category of embeddings, permitted manifolds, and surface conventions require clarification. The supplied record labels the problem open; that label is not independently verified here.
\subsection{Sources}
\begin{itemize}
\item Local UnsolvedMath registry: supplied complete record with permanent ID \(30004818\), exact source code \texttt{catalogue entry}, and title ``Genus-Reducing Local-Knot Concordance in Orientable Three-Manifolds.'' The record belongs to the Oberwolfach Reports Open Problems collection. No individual proposer is identified; the supplied proposed year is \(2021\).
\item Supplied report citation: \emph{MATRIX-MFO Tandem Workshop: Invariants and Structures in Low-Dimensional Topology} (2022), \emph{Oberwolfach Reports}, pp.~2327--2364, DOI \texttt{10.4171/owr/2021/42}. Source URL and literature-source link labelled ``OWR report'': \url{https://doi.org/10.4171/owr/2021/42}. This is supplied report-level provenance; the record gives no particular problem page or passage.
\item Supplied literature-source link labelled ``Celoria'': \url{https://doi.org/10.1112/topo.12051}. No title, full author information, or specific supporting result is supplied.
\item Literature assessment checked at \texttt{2026-08-22}: this date and assessment are supplied dated curation, not independently verified. The curation labels the orientable genus-saving question open and states that Boden--Nagel proves only the sliceness implication for local knots, not the requested genus separation. No bibliographic reference for Boden--Nagel is supplied. The background describes this assessment as dated literature triage derived from the supplied Oberwolfach curation and requests independent verification of time-sensitive status and cited claims.
\end{itemize}
\end{document}
